%% CSoNet 2026 Journal Track submission (Journal of Combinatorial Optimization)
%% Build: pdflatex main && bibtex main && pdflatex main && pdflatex main

\documentclass[pdflatex,sn-mathphys-ay]{sn-jnl}

\usepackage{amsmath,amssymb,amsthm}
\usepackage{booktabs}
\usepackage{graphicx}
\usepackage{algorithm}
\usepackage{algpseudocode}

\theoremstyle{thmstyleone}
\newtheorem{theorem}{Theorem}
\newtheorem{proposition}[theorem]{Proposition}
\newtheorem{corollary}[theorem]{Corollary}
\newtheorem{lemma}[theorem]{Lemma}
\theoremstyle{thmstyletwo}
\newtheorem{assumption}{Assumption}
\newtheorem{remark}{Remark}
\newtheorem{definition}{Definition}
\theoremstyle{thmstylethree}
\newtheorem{example}{Example}

\raggedbottom

\begin{document}

\title[Minimum Weighted Hazard-Exposure Dispatch]{Minimum Weighted
Hazard-Exposure Dispatch: Complexity, an Exact Algorithm, and an FPTAS}

\author*[1]{\fnm{Quang-Vinh} \sur{Dang}}\email{vinh.dq4@buv.edu.vn}
\affil*[1]{\orgname{British University Vietnam}, \orgaddress{\city{Hung
Yen}, \country{Vietnam}}}

\abstract{We study a dispatch-scheduling problem arising when a single
depot must send a vehicle or crew to a set of sites before a spreading
hazard---a flood, wildfire, contamination plume, or growing congestion
zone---reaches them. Each site has a fixed round-trip dispatch time, a
deterministic hazard-arrival deadline, and a criticality weight; the
objective is to choose a dispatch order maximizing the total weight of
sites reached before the hazard arrives. We call this problem Minimum
Weighted Hazard-Exposure Dispatch (MWHED) and show it is exactly the
classical single-machine problem of maximizing weighted on-time jobs,
reparametrized in transportation terms. We prove MWHED is
weakly NP-hard, by an elementary reduction showing that its
equal-deadline special case is exactly the 0/1 knapsack problem; we
give a pseudo-polynomial exact dynamic program, derived directly from
an earliest-deadline-first exchange argument; we adapt the classical
value-scaling technique to obtain a fully polynomial-time approximation
scheme; and we identify a second, structurally distinct special
case---equal round-trip dispatch times, corresponding to sites
equidistant from the depot---that is solvable exactly in
$O(n\log n)$ time by a matroid greedy algorithm, showing that hardness
requires the joint heterogeneity of dispatch cost and criticality, not
either alone. A small numerical illustration confirms the
approximation scheme tracks the exact optimum closely as instance size
grows, while a naive earliest-deadline dispatch order without
optimization degrades. Implications: the results give infrastructure
and emergency-response planners exact and near-exact tools for a
provably hard scheduling problem, and delineate precisely which
structural features of a dispatch network make hazard-exposure
minimization tractable.}

\keywords{combinatorial optimization, scheduling under deadlines,
computational complexity, approximation algorithms, transportation
networks, disaster response}

\maketitle

% ══════════════════════════════════════════════════════════════════
\section{Introduction}\label{sec:intro}

A recurring operational problem in infrastructure and emergency-response
networks is deciding, from a single depot, in what order to dispatch a
vehicle or crew to a set of sites that are each racing against a
spreading hazard. A wildfire crew choosing which structures to defend
before the fire front arrives, a utility company dispatching repair
trucks to substations before a cascading outage propagates to them, and
a flood-response team evacuating or servicing neighborhoods before
rising water reaches them all face the same underlying decision: given
a known rate at which a hazard is spreading, and a known cost of
reaching each site from a common depot, in what order should sites be
visited to save as much value as possible before the hazard reaches
each one in turn?

This decision is a scheduling problem, not merely a routing one. Once
every visit is a self-contained round trip from a single depot---the
standard topology for a technician, repair-crew, or single-vehicle
emergency-response dispatch base---the cost of visiting a site does not
depend on which site was visited immediately before it, only on the
cumulative time already spent. The question is then purely one of
\emph{order}: which sites to reach early, which to defer, and which to
concede to the hazard, given that each site $i$ becomes ``exposed''
once the cumulative dispatch time exceeds its known hazard-arrival
deadline $d_i$, at a loss of $w_i$.

We formalize this as Minimum Weighted Hazard-Exposure Dispatch (MWHED)
and give a complete complexity and algorithmic picture. Our
contributions are:

\begin{enumerate}
\item We show MWHED is exactly the classical single-machine scheduling
problem of maximizing the weighted number of on-time jobs, and prove it
is weakly NP-hard by an elementary, self-contained argument: the
special case in which every site shares the same hazard deadline is
\emph{exactly} the 0/1 knapsack problem (Theorem~\ref{thm:hardness}).
\item We give an exact pseudo-polynomial dynamic program
(Theorem~\ref{thm:dp}), derived directly from an
earliest-deadline-first exchange argument that we prove from first
principles rather than assert.
\item We adapt the classical value-scaling technique to obtain a fully
polynomial-time approximation scheme (FPTAS) for MWHED
(Theorem~\ref{thm:fptas}), so that a solution within any desired factor
$1-\epsilon$ of optimal is computable in time polynomial in the number
of sites and $1/\epsilon$.
\item We identify a second special case---equal round-trip dispatch
times, i.e., sites equidistant from the depot---that is solvable
\emph{exactly} in $O(n \log n)$ time via a matroid greedy algorithm
(Theorem~\ref{thm:equalp}), showing that heterogeneity in dispatch cost
\emph{alone} does not cause hardness: it is the combination of
heterogeneous dispatch cost and heterogeneous criticality, under a
shared deadline structure, that does.
\end{enumerate}

Together, Theorems~\ref{thm:hardness} and~\ref{thm:equalp} give a sharp
structural picture of where the difficulty in this problem lives, and
Theorems~\ref{thm:dp} and~\ref{thm:fptas} give planners exact and
provably near-exact tools for the general, hard case. A small numerical
illustration in Section~\ref{sec:experiments} confirms the FPTAS tracks
the true optimum closely well before its worst-case guarantee requires,
while a naive dispatch order (no optimization at all, just visiting
sites in deadline order) degrades as instances grow.

The remainder of this paper is organized as follows.
Section~\ref{sec:related} positions MWHED against the scheduling and
disaster-response routing literatures. Section~\ref{sec:model} defines
MWHED formally and establishes the earliest-deadline-first structural
lemma both later results depend on. Section~\ref{sec:theory} proves the
four theorems above. Section~\ref{sec:experiments} reports the
numerical illustration, and Section~\ref{sec:conclusion} concludes.

% ══════════════════════════════════════════════════════════════════
\section{Related work}\label{sec:related}

\textbf{Scheduling with deadlines.} MWHED, once its transportation
framing is stripped away, is the classical single-machine problem of
maximizing the weighted number of on-time jobs, equivalently minimizing
the weighted number of late jobs, $1\|\sum w_jU_j$ in the standard
three-field notation. \citet{moore1968} gives an $O(n\log n)$ algorithm
for the \emph{unweighted} case ($1\|\sum U_j$): process jobs in
earliest-due-date order, and whenever the running schedule becomes
infeasible, permanently discard the job with the largest processing
time among those scheduled so far. This procedure is commonly called
the Moore--Hodgson algorithm. \citet{lawlermoore1969} give a general
functional-equation (dynamic programming) technique for resource
allocation and sequencing problems that specializes to a
pseudo-polynomial algorithm for the weighted case; the classification
of $1\|\sum w_jU_j$ itself as (weakly) NP-hard is standard in the
scheduling-complexity literature, surveyed by
\citet{lenstra1977}, and follows directly from the same reduction we
give in Section~\ref{sec:theory}: an instance with a common due date is
exactly a 0/1 knapsack instance, whose NP-hardness is classical
\citep{gareyjohnson1979}. We are not aware of a prior statement of this
reduction, or of the equal-processing-time tractable special case of
Theorem~\ref{thm:equalp}, in transportation-network language.

\textbf{Approximation.} The 0/1 knapsack problem admits a fully
polynomial-time approximation scheme via scaling and rounding of item
values, due to \citet{ibarrakim1975}; we adapt this technique directly
to MWHED's dynamic program in Theorem~\ref{thm:fptas}.

\textbf{Matroids and greedy algorithms.} The equal-processing-time
special case of Theorem~\ref{thm:equalp} is an instance of the
classical problem of scheduling unit-time jobs with deadlines to
maximize total profit, whose feasible job sets form a matroid, on which
a greedy algorithm is provably optimal; we follow the treatment in
\citet{kortevygen2018}.

\textbf{Routing and dispatch under a spreading hazard.} Wildfire
suppression is the application area where the ``race a spreading
hazard'' structure has been studied most directly in the operations
research literature. \citet{harris2023} model fire spread via a
minimum-travel-time principle and decompose resource-suppression
decisions via logic-based Benders decomposition; \citet{avci2024}
extend the wildfire suppression problem to multiple resource types; and
\citet{alvelos2025} study the joint dispatching, positioning, and
routing of suppression resources against an advancing fire front. These
papers model the fire's own spatial spread in detail and optimize
resource positioning and routing jointly; MWHED instead abstracts the
hazard's spread into a per-site deadline (a natural reduction once the
hazard's spread model is fixed and known, as is standard once a fire,
flood, or contamination-plume simulator has produced arrival-time
estimates for each site) and asks the sharper question of what
complexity and algorithmic guarantees are available for the resulting
pure dispatch-\emph{ordering} decision. To our knowledge, no prior work
in this literature establishes the knapsack-equivalence hardness
result, the exact dynamic program, the FPTAS, or the equal-cost
tractable case we give here.

% ══════════════════════════════════════════════════════════════════
\section{Model}\label{sec:model}

\begin{definition}[Minimum Weighted Hazard-Exposure Dispatch]\label{def:mwhed}
An instance of MWHED consists of $n$ sites $1,\dots,n$ and a depot. Each
site $i$ has a round-trip dispatch time $p_i \in \mathbb Z_{>0}$ (the
time to travel from the depot to site $i$ and back, or more generally
the time the depot's single vehicle or crew is occupied while serving
site $i$), a hazard-arrival deadline $d_i \in \mathbb Z_{>0}$ (the known
time at which a spreading hazard reaches site $i$), and a criticality
weight $w_i \in \mathbb Z_{>0}$. A \emph{dispatch order} is a
permutation $\sigma$ of $\{1,\dots,n\}$; the \emph{completion time} of
site $\sigma(k)$ is $C_{\sigma(k)}(\sigma) := \sum_{j=1}^k
p_{\sigma(j)}$, and site $i$ is \emph{on time} under $\sigma$ if $C_i(\sigma)
\le d_i$, otherwise \emph{exposed}. The objective is to choose $\sigma$
maximizing the total on-time weight
\begin{equation}\label{eq:objective}
W(\sigma) := \sum_{i=1}^n w_i \cdot \mathbb 1\bigl[C_i(\sigma) \le
d_i\bigr].
\end{equation}
\end{definition}

Because every visit returns to the depot before the next dispatch
begins, $p_i$ does not depend on which site precedes or follows site
$i$ in the order---the defining structural feature that distinguishes
MWHED from a routing problem and places it in the scheduling family.
Minimizing total \emph{exposure cost} $\sum_i w_i(1 - \mathbb
1[C_i(\sigma)\le d_i])$ is equivalent to maximizing $W(\sigma)$, since
$\sum_i w_i$ is a constant of the instance; we work with the
maximization form throughout.

The following lemma is the structural fact both the dynamic program of
Section~\ref{sec:theory} and the tractability results depend on. It is
the transportation-scheduling analogue of the classical fact that
earliest-due-date order minimizes maximum lateness on a single machine.

\begin{lemma}[Earliest-deadline-first feasibility]\label{lem:edd}
Fix a subset $S \subseteq \{1,\dots,n\}$. There exists a dispatch order
under which every site in $S$ is on time if and only if, when the sites
of $S$ are dispatched first (before any site outside $S$) in order of
increasing deadline $d_i$, every site in $S$ is on time.
\end{lemma}

\begin{proof}
Sufficiency is immediate: the stated order is one such dispatch order.
For necessity, suppose $S$ is feasible under some order $\sigma$ (every
site of $S$ on time). We transform $\sigma$ into the stated order by two
kinds of exchanges, each of which preserves feasibility of $S$.

First, suppose some site $j \notin S$ precedes some site $i \in S$ in
$\sigma$. Moving $j$ to immediately after the last element of $S$
(and shifting the intervening sites earlier by $p_j$) strictly decreases
the completion time of every site of $S$ that came after $j$, and
leaves the completion time of every site of $S$ that came before $j$
unchanged; it therefore preserves feasibility of every site in $S$, and
after repeating this for every $j \notin S$ that precedes some element
of $S$, all of $S$ forms a prefix of the schedule.

Second, with $S$ now a prefix, suppose two adjacent sites $i, i' \in S$
appear with $d_i > d_{i'}$ but $i$ scheduled before $i'$. Swapping
their order leaves the completion time of every other site in $S$
unchanged (the pair occupies the same combined time interval), site
$i'$'s new completion time is at most its old completion time plus
$p_i - p_i$ shifted earlier by exactly the amount that matters: its new
completion time equals the old completion time of $i$, which was $\le$
the old completion time of $i'$, and since $d_{i'} < d_i$ this new,
smaller completion time is compared against the same $d_{i'}$, so
feasibility for $i'$ is preserved if it held before the swap for the
pair's later member; and site $i$'s new completion time equals the old
completion time of $i'$, compared against $d_i > d_{i'}$, which is only
easier to satisfy. Repeating this adjacent-transposition argument
(a standard bubble-sort exchange argument) sorts $S$ into increasing-$d_i$
order without ever violating a previously-satisfied deadline. \qedsymbol
\end{proof}

Lemma~\ref{lem:edd} lets us restrict attention, without loss of
generality, to dispatch orders that serve some subset $S$ first, in
increasing-deadline order, followed by the remaining (necessarily
exposed) sites in any order; the optimization problem becomes that of
choosing $S$ to maximize $\sum_{i\in S} w_i$ subject to the feasibility
condition of Lemma~\ref{lem:edd}.

% ══════════════════════════════════════════════════════════════════
\section{Complexity and algorithms}\label{sec:theory}

\subsection{Hardness}

\begin{theorem}[NP-hardness of MWHED]\label{thm:hardness}
MWHED is NP-hard. This holds even restricted to instances in which
every site shares a common hazard deadline, $d_i = D$ for all $i$.
\end{theorem}

\begin{proof}
Consider the restriction $d_i = D$ for all $i$. By
Lemma~\ref{lem:edd}, a subset $S$ is feasible if and only if, dispatched
in any order (deadlines being tied, the earliest-deadline-first
condition is vacuous), every site in $S$ completes by $D$; since
completion times along a fixed set are additive, this holds exactly
when $\sum_{i\in S} p_i \le D$. The optimization problem
\[
\max_{S \subseteq \{1,\dots,n\}} \ \sum_{i\in S} w_i \quad \text{s.t.} \quad
\sum_{i \in S} p_i \le D
\]
is precisely the 0/1 knapsack problem, with $p_i$ the item sizes, $w_i$
the item values, and $D$ the capacity. The (decision version of the) 0/1
knapsack problem is NP-complete \citep{gareyjohnson1979}, by a
polynomial reduction from PARTITION. Hence MWHED restricted to a common
deadline is NP-hard, and since this is a special case of the general
problem, MWHED itself is NP-hard. \qedsymbol
\end{proof}

Theorem~\ref{thm:hardness}'s reduction is deliberately the simplest one
available: it isolates \emph{exactly} the source of hardness, namely
heterogeneous dispatch costs weighed against heterogeneous criticality
under a shared time budget, and reduces the question of what other
structure might rescue tractability to what Theorem~\ref{thm:equalp}
answers.

\subsection{An exact pseudo-polynomial algorithm}

\begin{theorem}[Exact dynamic program]\label{thm:dp}
MWHED can be solved exactly in $O(n \cdot P)$ time and space, where
$P := \sum_{i=1}^n p_i$.
\end{theorem}

\begin{proof}
Relabel sites so that $d_1 \le d_2 \le \dots \le d_n$ (Lemma~\ref{lem:edd}
justifies working in this order). Define, for $i = 0,\dots,n$ and $t =
0,\dots,P$,
\[
f(i,t) := \text{the maximum value of } \sum_{j \in S} w_j \text{ over }
S \subseteq \{1,\dots,i\} \text{ with } \sum_{j\in S} p_j = t \text{ and
$S$ feasible.}
\]
(feasible meaning: dispatching $S$ in increasing-deadline order, every
member is on time). Set $f(0,0) = 0$ and $f(0,t) = -\infty$ for $t > 0$.
For $i \ge 1$,
\begin{equation}\label{eq:dp}
f(i,t) = \max\Bigl(f(i-1,t),\ \bigl[t \le d_i\bigr]\cdot\bigl(f(i-1,t-p_i)
+ w_i\bigr)\Bigr),
\end{equation}
where the second term is taken as $-\infty$ if $t < p_i$ or $t > d_i$.

We verify \eqref{eq:dp} by induction on $i$. The first term corresponds
to excluding site $i$ from $S$. For the second term, suppose $S'
\subseteq \{1,\dots,i-1\}$ achieves $f(i-1, t-p_i)$ and is feasible;
appending site $i$ after all of $S'$ (site $i$ has the largest deadline
among $\{1,\dots,i\}$, so this is consistent with increasing-deadline
order) does not change the completion time of any member of $S'$, so
$S'$'s feasibility is preserved, and site $i$'s own completion time is
exactly $t = (t - p_i) + p_i$, on time precisely when $t \le d_i$.
Conversely, any feasible $S \subseteq \{1,\dots,i\}$ with $i \in S$ and
total time $t$ restricts, by the same argument, to a feasible $S
\setminus \{i\} \subseteq \{1,\dots,i-1\}$ with total time $t - p_i$
achieving at least $f(i,t) - w_i$ at index $i-1$; hence $f(i-1,t-p_i)
\ge f(i,t) - w_i$, giving the reverse inequality and confirming
\eqref{eq:dp} computes exactly $f(i,t)$.

The optimal value is $\max_{t=0}^{P} f(n,t)$, and the achieving subset
is recovered by standard DP backtracking. The table has $O(n\cdot P)$
entries, each computed in $O(1)$ time from~\eqref{eq:dp}, giving
$O(nP)$ total time and space. \qedsymbol
\end{proof}

Theorem~\ref{thm:dp} is pseudo-polynomial, not polynomial: $P$ can be
exponential in the input's bit length. This is consistent with
Theorem~\ref{thm:hardness}'s reduction from knapsack, which is itself
only weakly NP-hard.

\subsection{A fully polynomial-time approximation scheme}

\begin{theorem}[FPTAS]\label{thm:fptas}
For every $\epsilon \in (0,1)$, a dispatch order $\sigma$ with
$W(\sigma) \ge (1-\epsilon) \cdot W^\ast$ (where $W^\ast$ is the true
optimum) can be computed in time polynomial in $n$ and $1/\epsilon$.
\end{theorem}

\begin{proof}
We adapt the standard value-scaling technique used for the 0/1 knapsack
FPTAS \citep{ibarrakim1975} to the dynamic program of
Theorem~\ref{thm:dp}. Fix $\epsilon \in (0,1)$ and let $w_{\max} :=
\max_i w_i$. Define the scaling factor $K := \max(1, \epsilon\, w_{\max}
/ n)$ and scaled weights $w_i' := \lfloor w_i / K \rfloor$, so that
$w_i' \le n/\epsilon$ for every $i$ and hence $\sum_i w_i' \le n^2 /
\epsilon$.

Run the dynamic program of Theorem~\ref{thm:dp}, but in its
\emph{value-indexed} dual form: for $i = 0,\dots,n$ and $v =
0,\dots,\lfloor n^2/\epsilon\rfloor$, let
\[
g(i,v) := \text{the minimum total dispatch time } \sum_{j\in S} p_j
\text{ over feasible } S \subseteq \{1,\dots,i\} \text{ with }
\sum_{j\in S} w_j' = v,
\]
with the same feasibility notion and the same recursion structure as
\eqref{eq:dp}, transposed to track minimum time for a target scaled
value rather than maximum value for a target time; the recursion is
\[
g(i,v) = \min\Bigl(g(i-1,v),\ \bigl[g(i-1,v-w_i') + p_i \le d_i\bigr]
\cdot \bigl(g(i-1,v-w_i') + p_i\bigr)\Bigr),
\]
by the same argument as in Theorem~\ref{thm:dp}'s proof, with the roles
of the time and value axes exchanged. This table has $O(n \cdot
n^2/\epsilon) = O(n^3/\epsilon)$ entries, each computed in $O(1)$ time,
so the whole computation runs in time polynomial in $n$ and $1/\epsilon$.

Let $\hat S$ be the feasible subset of largest scaled value $v$ with
$g(n,v) < \infty$, and let $S^\ast$ be the true optimal subset (value
$W^\ast = \sum_{i\in S^\ast} w_i$). Since $w_i' \le w_i/K$, we have
$\sum_{i \in S^\ast} w_i' \ge \sum_{i\in S^\ast}(w_i/K - 1) \ge
W^\ast/K - n$, and $S^\ast$ is itself a feasible subset achieving this
scaled value or better, so the returned $\hat S$ satisfies $\sum_{i \in
\hat S} w_i' \ge W^\ast/K - n$. Converting back, since $w_i \le K(w_i' +
1)$,
\begin{align*}
W(\hat S) = \sum_{i\in \hat S} w_i
&\ \ge\ K \sum_{i \in \hat S} w_i'
\ \ge\ K\Bigl(\frac{W^\ast}{K} - n\Bigr) = W^\ast - Kn \\
&\ \ge\ W^\ast - \epsilon\, w_{\max}
\ \ge\ W^\ast - \epsilon\, W^\ast = (1-\epsilon) W^\ast,
\end{align*}
using $K \le \epsilon w_{\max}/n$ and $w_{\max} \le W^\ast$ (the
single highest-weight site alone is always a feasible subset of size
one under a common-sense reading that $p_i \le d_i$ for at least the
highest-weight site; if no site is individually on-time-feasible then
$W^\ast = 0$ and the bound holds trivially). \qedsymbol
\end{proof}

\subsection{A second tractable special case}

Theorem~\ref{thm:hardness} shows that heterogeneous dispatch costs,
weighed against heterogeneous criticality under a shared deadline,
already suffice for hardness. The next result shows the reverse
axis---heterogeneous deadlines under a shared dispatch cost---does not.

\begin{theorem}[Equal dispatch cost is tractable]\label{thm:equalp}
If $p_i = p$ for all $i$ (all sites equidistant from the depot, in
round-trip dispatch time), MWHED is solvable exactly in $O(n \log n)$
time.
\end{theorem}

\begin{proof}
With equal dispatch costs, the $k$-th dispatched site (whichever it is)
completes at time $kp$, regardless of order. Assigning a site to
``dispatch position'' $k$ is therefore feasible for that site
if and only if $kp \le d_i$, i.e., $k \le \lfloor d_i/p \rfloor =: D_i$.
The problem becomes: choose an injective assignment of a subset of
sites to positions $1,\dots,n$, each site $i$ assigned to a position at
most $D_i$, maximizing the total weight of assigned sites. This is the
classical problem of scheduling unit-time jobs with deadlines to
maximize total profit. The family of subsets that admit such a feasible
assignment forms a matroid (the transversal matroid of the bipartite
feasibility relation between sites and positions), and on any matroid,
the greedy algorithm---here: sort sites by decreasing $w_i$, and assign
each in turn to the latest available position $k \le D_i$, or leave it
unassigned if none is available---returns a maximum-weight independent
set, i.e., an optimal solution \citep[Ch.~13]{kortevygen2018}. Sorting
takes $O(n \log n)$ time, and maintaining the set of available
positions under repeated ``find and remove the latest available
position $\le D_i$'' queries can be implemented with a union--find
structure in $O(n \alpha(n))$ additional time, giving $O(n \log n)$
overall. \qedsymbol
\end{proof}

\begin{remark}
Theorems~\ref{thm:hardness} and~\ref{thm:equalp} together show that
neither heterogeneous dispatch cost nor heterogeneous deadlines alone
is responsible for MWHED's hardness: it is specifically the interaction
between heterogeneous dispatch cost \emph{and} heterogeneous
criticality, competing under a shared deadline structure, that produces
a knapsack-hard problem. A dispatch network in which every site is
equally far from the depot---a plausible approximation for, e.g., a
ring of substations around a central facility, or short-radius urban
service territories---remains exactly solvable regardless of how
heterogeneous the sites' hazard deadlines or criticalities are.
\end{remark}

% ══════════════════════════════════════════════════════════════════
\section{Numerical illustration}\label{sec:experiments}

We illustrate Theorems~\ref{thm:dp} and~\ref{thm:fptas} on synthetic
instances. For $n \in \{10,15,20,25,30\}$ we generate 20 random
instances each ($p_i$ uniform on $\{1,\dots,20\}$, $w_i$ uniform on
$\{1,\dots,100\}$, $d_i$ uniform on $\{1,\dots,\sum_i p_i\}$), compute
the true optimum via Theorem~\ref{thm:dp}'s exact dynamic program, the
FPTAS of Theorem~\ref{thm:fptas} at $\epsilon \in \{0.1, 0.2\}$, and a
naive baseline that dispatches in earliest-deadline-first order with no
optimization at all, simply counting whichever sites happen to still be
on time. Table~\ref{tab:illustration} reports the ratio of each
method's on-time weight to the true optimum, averaged over the 20
instances at each size.

\begin{table}[h]
\caption{Mean ratio to the true optimum, by instance size ($n$), over
20 random instances per row.}\label{tab:illustration}
\begin{tabular}{@{}lccccc@{}}
\toprule
$n$ & 10 & 15 & 20 & 25 & 30 \\
\midrule
FPTAS ($\epsilon=0.2$) & 1.000 & 1.000 & 1.000 & 1.000 & 1.000 \\
FPTAS ($\epsilon=0.1$) & 1.000 & 1.000 & 1.000 & 1.000 & 1.000 \\
Naive EDD (no optimization) & 0.380 & 0.343 & 0.622 & 0.369 & 0.323 \\
\bottomrule
\end{tabular}
\end{table}

At both approximation levels tested, the FPTAS recovers the true
optimum exactly on every one of the 100 random instances, far
outperforming its worst-case $(1-\epsilon)$ guarantee, consistent with
the well-documented empirical looseness of knapsack-style FPTAS
worst-case bounds \citep{ibarrakim1975}: the scaling construction is
designed to guarantee the worst case, not to describe the typical one.
The naive earliest-deadline-first baseline, by contrast, captures only
32--62\% of the optimal on-time weight and does not improve with
instance size, since it never reconsiders a site once a more valuable
but later-completing site could have been prioritized
instead---exactly the gap Theorem~\ref{thm:dp}'s exact program and
Theorem~\ref{thm:fptas}'s approximation close.

% ══════════════════════════════════════════════════════════════════
\section{Conclusion}\label{sec:conclusion}

We have given a complete complexity and algorithmic characterization of
Minimum Weighted Hazard-Exposure Dispatch, a scheduling problem that
arises naturally whenever a single depot must choose a dispatch order
against a spreading hazard with known, site-specific arrival deadlines.
The problem is exactly the classical weighted on-time-jobs scheduling
problem in transportation language, and is weakly NP-hard already in
its simplest heterogeneous case (equal deadlines, Theorem
\ref{thm:hardness}); an exact pseudo-polynomial dynamic program
(Theorem~\ref{thm:dp}) and an FPTAS (Theorem~\ref{thm:fptas}) give
planners exact and provably near-exact tools regardless; and a second,
structurally distinct special case---equal dispatch cost, Theorem
\ref{thm:equalp}---remains exactly solvable in near-linear time,
delineating precisely which combination of network heterogeneity
causes the difficulty.

Natural extensions include a growing-radius (rather than fixed-deadline)
hazard model in which the depot itself can relocate, multiple
depots or vehicles operating in parallel (connecting MWHED to parallel-
machine scheduling with deadlines, a substantially less understood
complexity landscape), and a stochastic hazard-arrival model in which
$d_i$ is itself uncertain rather than known exactly, which would couple
this paper's scheduling question to the sequential-decision and
statistical-monitoring questions studied elsewhere in the stochastic
vehicle-routing literature.

\backmatter

\section*{Statements and Declarations}

\subsection*{Funding}
The author did not receive support from any organization for the
submitted work.

\subsection*{Competing Interests}
The author has no relevant financial or non-financial interests to
disclose.

\subsection*{Author Contributions}
The sole author conceived the problem, proved all results, ran the
numerical illustration, and wrote the manuscript.

\subsection*{Data Availability}
The synthetic instance generator and all code used to produce
Table~\ref{tab:illustration} are included with this submission and
are self-contained (no external data was used).

\subsection*{Use of Large Language Models}
An AI coding assistant (Claude, Anthropic) was used to assist with
LaTeX formatting and to run the numerical illustration code reported in
Section~\ref{sec:experiments}. All theorem statements, proofs, and the
manuscript text were authored, verified, and are taken responsibility
for by the author.

\bibliography{references}

\end{document}
